\documentclass[11pt,letterpaper]{article}
\usepackage[T1]{fontenc}
\usepackage{lmodern}
\usepackage[margin=1in]{geometry}
\usepackage{amsmath,amssymb,amsthm,mathtools}
\usepackage{microtype}
\usepackage{enumitem,longtable,booktabs,array}
\usepackage[colorlinks=true,linkcolor=black,citecolor=black,urlcolor=blue]{hyperref}
\usepackage{bookmark}
\hypersetup{pdftitle={An explicit invariant for a degree five symbol descent family},pdfsubject={Scoped degree-five cyclicity with checked Lean proof and staged replay; original earlier object reuse disclosed}}
\setlength{\parskip}{0.2em}
\setlength{\emergencystretch}{2em}
\allowdisplaybreaks[2]
\newtheorem{theorem}{Theorem}[section]
\newtheorem{proposition}[theorem]{Proposition}
\newtheorem{lemma}[theorem]{Lemma}
\newtheorem{corollary}[theorem]{Corollary}
\theoremstyle{definition}
\newtheorem{definition}[theorem]{Definition}
\theoremstyle{remark}
\newtheorem{remark}[theorem]{Remark}
\newcommand{\ZZ}{\mathbb Z}
\newcommand{\QQ}{\mathbb Q}
\newcommand{\FF}{\mathbb F}
\newcommand{\Gal}{\operatorname{Gal}}
\newcommand{\supp}{\operatorname{supp}}
\newcommand{\ev}{\operatorname{ev}}
\newcommand{\Cent}{\operatorname{Z}}
\newcommand{\Tr}{\operatorname{Tr}}
\newcommand{\id}{\operatorname{id}}
\newcommand{\red}{\operatorname{red}}
\newcommand{\code}[1]{\texttt{\detokenize{#1}}}
\title{An explicit invariant for a degree five symbol descent family}
\author{}
\date{October 13, 2026}
\begin{document}
\maketitle
\begin{abstract}
Let $K$ be a characteristic-zero field containing a primitive fifth root of
unity $q$. Consider a central division algebra with symbol generators
$U,V$, relation $UV=qVU$, scalar fifth powers, and the actual $25$-element
symbol basis. Suppose it admits an order-three semilinear automorphism
cycling $U,V,(UV)^{-1}$ and fixing $q$. We describe a $69$-term Laurent
expression $N$, with a short expression using ordered pairs, which is
fixed and nonscalar in every such division-algebra specialization. An
elementary interpolation argument proves that the fixed division algebra
has degree five over the fixed scalar field. An exact finite identity
$N^5=H(U^5,V^5)$, with a displayed $75$-term integer Laurent polynomial $H$,
gives a cyclic maximal subfield generated by $N$. The construction
is motivated by Fryer's invariant generators and uses no primitive cube
root in the hypotheses. The complete fifth-power certificate and the
relative-family theorem have passed Lean checking, with no extra power
hypothesis in the final theorem. A staged replay has also passed, with
earlier checked objects reused as detailed below. Historical novelty is
unestablished. The argument concerns this explicit
semilinear symbol family and does not settle unrestricted prime-degree
cyclicity.
\end{abstract}

\paragraph{Verification boundary.}
The full evaluated fifth-power certificate and the final relative-family
assembly have passed Lean checking. The final export assumes exactly the
structural symbol data stated below, with no additional numerator or
power hypothesis. Its audit reports only the standard logical principles
\code{propext}, \code{Classical.choice}, and \code{Quot.sound}; the literal
output alignment uses no axioms. The N4 arithmetic and structural
rebuilds, fresh N5 checks, and downstream final-assembly checks have
passed. The final staged replay reuses $99$ original earlier base/N2/N3
objects; it is not a fresh rebuild of the entire arithmetic closure.
The distinction between these checks, the mathematical argument, and the
additional paper-level description of a Galois generator is recorded in
Section~\ref{sec:formal}. No novelty or priority claim is made.

\section{The precise family and the conclusion}\label{sec:setup}

All division algebras and field homomorphisms are unital. We identify a field
with its specified central image when this causes no ambiguity. Negative
powers are powers of inverses. Products are written in their intended order;
coefficients in the center may be moved past all factors, but arbitrary
algebra elements may not.

\begin{definition}[Symbol data with an order-three action]\label{def:data}
The data consist of the following objects and hypotheses.
\begin{enumerate}[label=\textup{(\roman*)},leftmargin=2.4em]
\item A field $K$ of characteristic zero, a division algebra $A$ with a
central embedding $K\hookrightarrow A$, and $\Cent(A)=K$.
\item An element $q\in K$ of multiplicative order five and units $U,V\in A$
satisfying
\begin{equation}\label{eq:quantum}
UV=qVU.
\end{equation}
\item Elements $x,y\in K$ such that $U^5=x$ and $V^5=y$, and the assertion
that
\begin{equation}\label{eq:basis}
\{U^iV^j:0\leq i,j<5\}
\end{equation}
is a $K$-basis of $A$.
\item Ring automorphisms $\sigma\colon A\to A$ and $\tau\colon K\to K$
satisfying, for every $c\in K$ and $a\in A$,
\begin{equation}\label{eq:action}
\begin{gathered}
\sigma(ca)=\tau(c)\sigma(a),\qquad \sigma^3=\id,\\
\sigma(U)=V,\qquad \sigma(V)=(UV)^{-1},\qquad \tau(q)=q.
\end{gathered}
\end{equation}
\end{enumerate}
Put
\begin{equation}\label{eq:fixed}
W=(UV)^{-1},\qquad z=(xy)^{-1},\qquad
F=K^\tau,\qquad D=A^\sigma.
\end{equation}
\end{definition}

The fifth powers $x,y$ are nonzero because $U,V$ are units. The basis
hypothesis implies $\dim_K A=25$. We will prove, rather than assume,
$\dim_F D=25$ and $\Cent(D)=F$. Likewise, invariance, nonscalarity, and the
central fifth power of the particular numerator are not part of the data.
Although $\sigma^3=\id$ allows order dividing three in its wording, the
hypotheses force both scalar and algebra automorphisms to have exact order
three; see Lemma~\ref{lem:cycle} and Lemma~\ref{lem:distinct}.

For an intrinsic central division algebra with a ring automorphism
$\sigma$, restriction to the center supplies $\tau$, and multiplicativity
supplies semilinearity. Definition~\ref{def:data} spells out these maps
because the formal interface does not identify a field literally with a
subset of a different type. It introduces no assumption about a general
Galois descent theorem.

\paragraph{Family conclusion.}
For every datum in Definition~\ref{def:data}, the explicitly specified
$N\in A$ below generates a cyclic maximal subfield $F(N)$ of $D/F$, of
degree five. The proof uses the exact Laurent identity in
Proposition~\ref{prop:fifth}. The full conclusion is stated in
Theorem~\ref{thm:main}.

This is a relative assertion about a supplied symbol presentation and a
particular action. It does not assert that every degree-five division
algebra admits such data, that every cubic descent of a cyclic algebra is
cyclic, or that arbitrary scalar choices $x,y$ give a division algebra.
No construction of a generic Ore division ring is needed anywhere in the
argument. Constructing a generic example is separate background and is not
part of the current formal theorem.

\section{Quantum monomials and the scalar cycle}\label{sec:monomials}

We first record all commutation conventions, since a phase error changes
the finite certificate.

\begin{lemma}[Integer exponent rules]\label{lem:monomial}
For integers $i,j,k,l$ one has
\begin{align}
V^jU^k&=q^{-jk}U^kV^j,\label{eq:commute}\\
(U^iV^j)(U^kV^l)&=q^{-jk}U^{i+k}V^{j+l},\label{eq:mul}\\
(U^iV^j)^n&=q^{-ij n(n-1)/2}U^{in}V^{jn}\qquad(n\geq0).
\label{eq:power}
\end{align}
In particular, $(UV)^5=xy$, $W^5=z$, and
\begin{equation}\label{eq:ratio-power}
(UV^{-1})^5=xy^{-1}.
\end{equation}
\end{lemma}
\begin{proof}
Equation~\eqref{eq:quantum} gives $VU=q^{-1}UV$. Multiplying this equality
on the appropriate side by inverses gives the corresponding relations with
$U^{-1}$ and $V^{-1}$. Induction first on a nonnegative exponent and then on
its negative establishes \eqref{eq:commute} for all integers. Multiplying
$U^i(V^jU^k)V^l$ gives \eqref{eq:mul}.

For \eqref{eq:power}, the assertion is immediate when $n=0$. Multiplication
on the right by $U^iV^j$ adds the phase $q^{-ij n}$, and
$n(n-1)/2+n=n(n+1)/2$. This proves the induction. With $n=5$ the phase is
$q^{-10ij}=1$, because $q^5=1$. Thus $(UV)^5=U^5V^5=xy$, and taking
inverses gives $W^5=z$. The same calculation with $(i,j)=(1,-1)$ gives
\eqref{eq:ratio-power}.
\end{proof}

\begin{lemma}[The induced scalar action]\label{lem:cycle}
The automorphism $\sigma$ cycles $U,V,W$, and $\tau$ cycles $x,y,z$.
Moreover $q\in F$, and it has order five in $F$.
\end{lemma}
\begin{proof}
Since $W=V^{-1}U^{-1}$, we have $VW=U^{-1}$. Therefore
$\sigma(W)=(\sigma(U)\sigma(V))^{-1}=(VW)^{-1}=U$.
Applying semilinearity to $c\cdot1$ shows $\sigma(c)=\tau(c)$ on the
specified copy of $K$. Consequently
\[
\tau(x)=\sigma(U^5)=V^5=y,\qquad
\tau(y)=\sigma(V^5)=W^5=z,\qquad
\tau(z)=\sigma(W^5)=U^5=x.
\]
Finally $\tau(q)=q$ places $q$ in $F$. The inclusion $F\hookrightarrow K$
is injective, so it preserves the least positive exponent giving $q^n=1$.
\end{proof}

\begin{lemma}[A centrally split polynomial in a division ring]\label{lem:split}
If $t\in A$ satisfies $t^5=1$, then $t=q^i$ for some $0\leq i<5$.
\end{lemma}
\begin{proof}
The elements $1,q,\ldots,q^4$ are distinct, and
$X^5-1=\prod_{i=0}^4(X-q^i)$ in $K[X]$. Evaluation at $t$ respects
multiplication because $t$ commutes with every coefficient. Thus
$0=\prod_{i=0}^4(t-q^i)$ in $A$. A division ring has no zero divisors:
if $ab=0$ and $a\ne0$, multiplication by $a^{-1}$ gives $b=0$.
Repeated use of this fact shows that one of the displayed factors is zero.
\end{proof}

\begin{lemma}[Pairwise distinct scalar parameters]\label{lem:distinct}
The elements $x,y,z$ are pairwise distinct. In particular,
$(x-y)(x-z)\ne0$.
\end{lemma}
\begin{proof}
Suppose $x=y$. By \eqref{eq:ratio-power}, $t=UV^{-1}$ has $t^5=1$.
Lemma~\ref{lem:split} makes $t$ central, so $U=q^iV$ for some $i$.
Then $U$ and $V$ commute. Comparing with \eqref{eq:quantum} yields
$(1-q)VU=0$, impossible since $VU\ne0$ and $q\ne1$. Hence $x\ne y$.
Applying the injective map $\tau$ to this inequality and then applying it
again gives $y\ne z$ and $z\ne x$. The nonzero product assertion follows
in the field $K$. Restricting $\sigma^3=\id$ to the embedded field gives
$\tau^3=\id$. Since $\tau(x)=y\ne x$, $\tau$ has exact order three,
and the same is then true of $\sigma$.
\end{proof}

\begin{remark}
This proof uses the division hypothesis essentially. A noncentral matrix
can satisfy $t^5=1$, so the same argument cannot be transferred to split
matrix algebras just because their formal symbol relations look similar.
The conclusions of this paper apply to division-algebra specializations
satisfying Definition~\ref{def:data}.
\end{remark}

\section{Elementary descent of dimension and center}\label{sec:descent}

The following interpolation calculation supplies exactly the descent
statements needed here. It avoids placing the desired fixed-algebra
properties in the hypotheses.

\begin{proposition}[A fixed spanning formula]\label{prop:span}
For $a\in A$, define
\begin{align*}
t_0(a)&=a+\sigma(a)+\sigma^2(a),\\
t_1(a)&=xa+y\sigma(a)+z\sigma^2(a),\\
t_2(a)&=x^2a+y^2\sigma(a)+z^2\sigma^2(a).
\end{align*}
Each $t_i(a)$ belongs to $D$, and
\begin{equation}\label{eq:interpolation}
a=\frac{t_2(a)-(y+z)t_1(a)+yz\,t_0(a)}{(x-y)(x-z)}.
\end{equation}
Consequently $D$ spans $A$ over $K$.
\end{proposition}
\begin{proof}
Applying $\sigma$ to $t_0(a)$ cyclically permutes its three summands.
For $t_1(a)$, semilinearity and Lemma~\ref{lem:cycle} give
\[
\sigma(t_1(a))=y\sigma(a)+z\sigma^2(a)+xa=t_1(a).
\]
The calculation for $t_2(a)$ is identical with squared coefficients.
Expand the numerator of \eqref{eq:interpolation}, retaining the order of
the algebra elements. Since $x,y,z$ are central, the coefficients of
$a,\sigma(a),\sigma^2(a)$ are respectively
\[
x^2-(y+z)x+yz=(x-y)(x-z),\quad
y^2-(y+z)y+yz=0,\quad
z^2-(y+z)z+yz=0.
\]
The denominator is a nonzero central scalar by
Lemma~\ref{lem:distinct}; division is therefore legitimate.
Equation~\eqref{eq:interpolation} expresses every $a$ as a $K$-linear
combination of three elements of $D$.
\end{proof}

\begin{theorem}[Structure of the fixed algebra]\label{thm:descent}
Under Definition~\ref{def:data}, $D$ is a central division algebra over $F$
with
\[
\dim_F D=25,\qquad \Cent(D)=F.
\]
Thus $D$ has degree five.
\end{theorem}
\begin{proof}
Fixed elements are closed under addition, subtraction, and multiplication.
If $d\in D$ is nonzero, then
$\sigma(d^{-1})=\sigma(d)^{-1}=d^{-1}$, so $D$ is a division ring.
Likewise $F$ is a field. Its elements lie in $D$ and are central there.

By Proposition~\ref{prop:span}, a $K$-basis of $A$ can be selected from
$D$: begin with the empty independent family and repeatedly select a fixed
element outside its $K$-span until the span is all of $A$. This terminates
with $25$ elements because \eqref{eq:basis} gives $\dim_K A=25$.
Write this selected basis as $d_1,\ldots,d_{25}$. It is $F$-linearly
independent, because any relation with coefficients in $F$ is also a
relation with coefficients in $K$.

For any $d\in D$ there are unique $k_i\in K$ with $d=\sum_i k_i d_i$.
Apply $\sigma$ and use fixedness of $d,d_i$:
\[
\sum_i k_i d_i=d=\sum_i\tau(k_i)d_i.
\]
Uniqueness of coordinates in a $K$-basis gives $\tau(k_i)=k_i$ for every
$i$. Hence all $k_i$ lie in $F$, and the selected family is an $F$-basis
of $D$. This proves the dimension assertion.

Now let $c\in\Cent(D)$. It commutes with each $t_i(a)$ in
\eqref{eq:interpolation} and with every scalar in $K$, because
$K=\Cent(A)$. That equation implies $ca=ac$ for all $a\in A$.
Therefore $c\in\Cent(A)=K$, and its fixedness gives $c\in F$.
The reverse inclusion $F\subseteq\Cent(D)$ was already noted.
A finite-dimensional central division algebra has degree equal to the
square root of its central dimension; here that degree is five.
\end{proof}

The same selected basis shows that the multiplication map
$D\otimes_F K\to A$, $d\otimes k\mapsto kd$, is an isomorphism of
$K$-algebras: it maps the $K$-basis $d_i\otimes1$ to the $K$-basis $d_i$.
This observation explains the term ``descent'', but neither a tensor-product
construction nor the general Galois descent theorem is used to prove
Theorem~\ref{thm:descent}.

\section{A compact numerator without cube roots}\label{sec:compact}

A short expression is useful for proving invariance. A separate expanded
expression is useful for exact multiplication and basis coordinates. We
specify both and identify the bridge between them.

\subsection{Ordered pairs and a formal central variable}

For pairs of elements of $A$, define
\begin{equation}\label{eq:pair-product}
(a,b)\star(d,e)=(ad-be,\;ae+bd-be).
\end{equation}
Addition and scalar multiplication are componentwise. This is the
multiplication obtained by writing $a+\xi b$ in the ring
\[
A[\xi]/(\xi^2+\xi+1),
\]
where $\xi$ is a formal central variable and the quotient is by the
two-sided ideal generated by the central monic polynomial
$\xi^2+\xi+1$. Monic division by
$\xi^2+\xi+1$ gives a unique representative of degree at most one; the
product is therefore associative even though $A$ need not be commutative.
This auxiliary ring can have zero divisors. We never assert that it is a
field or a division algebra, and we take no inverses in it.

In this pair notation, multiplication by $\xi^2$ and $\xi$ is given by
\begin{equation}\label{eq:TS}
T(a,b)=(b-a,-a),\qquad S(a,b)=(-b,a-b).
\end{equation}
These equalities follow directly from $\xi^2=-1-\xi$ and $\xi^3=1$.
Both operations commute with multiplication by central scalars. The map
$\sigma$ acts componentwise and respects $\star$, since the product only
uses sums, differences, and products in $A$.

\subsection{Definition of the numerator}

Set
\begin{align}
\mathcal A&=(U-W,\;V-W),&
\mathcal B&=(U^{-1}-W^{-1},\;V^{-1}-W^{-1}),\label{eq:AB}\\
c&=5(1+q^{-1}),&
(d_0,d_1)&=(x-y,\;z-y),\label{eq:cd}\\
(P_0,P_1)&=(\mathcal A^{\star4}-c\mathcal A)\star\mathcal B.
\label{eq:P}
\end{align}
Here $\mathcal A^{\star4}=(\mathcal A\star\mathcal A)\star
(\mathcal A\star\mathcal A)$, and $c$ multiplies both components on the
left. Define
\begin{equation}\label{eq:Ncompact}
N_{\mathrm{cmp}}=(2d_0-d_1)P_0+(2d_1-d_0)P_1.
\end{equation}
Every inverse in this expression is the inverse of one of the specified
units or of $q$. In particular, the formula has no denominator involving
$x-y$, a discriminant, or a sum of scalar monomials. It is defined for
every datum in Definition~\ref{def:data}.

\begin{proposition}[Invariance of the compact formula]\label{prop:invariance}
One has $\sigma(N_{\mathrm{cmp}})=N_{\mathrm{cmp}}$.
\end{proposition}
\begin{proof}
Using the cycle $U\mapsto V\mapsto W\mapsto U$,
\[
\sigma(\mathcal A)=(V-U,W-U)=T\mathcal A.
\]
The same calculation for the inverses gives
$\sigma(\mathcal B)=T\mathcal B$. In the formal-variable notation, these
are multiplication by $\xi^2$. Since $c$ is fixed, $\xi$ is central, and
$\xi^8=\xi^2$, we obtain
\[
\sigma(\mathcal A^{\star4}-c\mathcal A)
 =\xi^2(\mathcal A^{\star4}-c\mathcal A),\qquad
\sigma(P_0,P_1)=\xi^4(P_0,P_1)=S(P_0,P_1).
\]
This calculation involves no specialization of $\xi$ to an actual cube
root in $K$.

The scalar pair transforms in the same way:
\[
\tau(d_0,d_1)=(y-z,x-z)=(-d_1,d_0-d_1)=S(d_0,d_1).
\]
For a central pair $d=(d_0,d_1)$ and an arbitrary pair $P=(P_0,P_1)$,
write $L(d,P)=(2d_0-d_1)P_0+(2d_1-d_0)P_1$.
A direct calculation, keeping all coefficients on the left, gives
\begin{align*}
L(Sd,SP)
 &=(-d_0-d_1)(-P_1)+(2d_0-d_1)(P_0-P_1)\\
 &=(2d_0-d_1)P_0+(2d_1-d_0)P_1=L(d,P).
\end{align*}
Applying this identity to the two transformed pairs proves the assertion.
\end{proof}

\subsection{The expanded numerator and the checked bridge}

Let $\mathcal S\subset\ZZ^2$ be the $69$ exponent pairs in
Appendix~\ref{app:N}. For each $(i,j)\in\mathcal S$, the table gives an
integer quadruple $(a_0,a_1,a_2,a_3)$, defining
$n_{ij}(q)=a_0+a_1q+a_2q^2+a_3q^3$. Put
\begin{equation}\label{eq:Nexpanded}
N=\sum_{(i,j)\in\mathcal S}n_{ij}(q)U^iV^j.
\end{equation}

\begin{proposition}[Identification of the two formulas]\label{prop:bridge}
With the phase convention \eqref{eq:mul}, one has $N=N_{\mathrm{cmp}}$.
Consequently $N\in D$.
\end{proposition}
\begin{proof}[Finite expansion and its verification]
Replace $W$ by $q^{-1}U^{-1}V^{-1}$ and $W^{-1}$ by $UV$ in
\eqref{eq:AB}; the first replacement follows from
$W=V^{-1}U^{-1}$ and \eqref{eq:commute}. Replace $x,y,z$ by
$U^5,V^5,U^{-5}V^{-5}$ in \eqref{eq:cd}. Expand each of the pair products
using \eqref{eq:pair-product}, multiply ordered monomials using
\eqref{eq:mul}, and reduce coefficients with
\begin{equation}\label{eq:cyclotomic}
q^4=-1-q-q^2-q^3.
\end{equation}
The resulting coefficients are exactly those in Appendix~\ref{app:N}.
This is a finite coefficient assertion, rather than a commutative
simplification of the noncommuting variables. The accompanying replay
script explicitly performs these replacements and compares all coefficients,
including absent coefficients, with the table. The recorded formal
compact-to-Laurent bridge also proves this equality after evaluation in
the actual division algebra; see Section~\ref{sec:formal}. Thus invariance
from Proposition~\ref{prop:invariance} applies to the specific expanded
numerator, not merely to a different expression with the same name.
\end{proof}

\paragraph{Relation to the motivating construction.}
Fryer's order-three construction uses an auxiliary primitive cube root
$\omega$ and an invariant expressed as $a^{-1}b$; see
\cite[Proposition~3.3 and equations (3.2.1)--(3.2.2)]{Fryer}.
Our pairs encode the expressions
$a=U+\omega V+\omega^2W$ and
$b=U^{-1}+\omega V^{-1}+\omega^2W^{-1}$ without requiring that root in $K$.
For a pair $p=p_0+\xi p_1$, the formal substitution
$\xi\mapsto\xi^2$ gives $\overline p=(p_0-p_1)-\xi p_1$.
The sum $p+\overline p$ is $2p_0-p_1$.
In particular, \eqref{eq:Ncompact} is the formal trace of
$(P_0+\xi P_1)\overline{(d_0+\xi d_1)}$.
This describes the trace-like shape of the formula; the proof uses the
explicit polynomial expression, not an assertion that inverses survive an
extension or specialization. No denominator-clearing step is used to
infer nonscalarity after specialization.

\section{Uniform nonscalarity from two coordinates}\label{sec:nonscalar}

A noncentral term in a generic Laurent polynomial need not remain a
nonzero noncentral coordinate after specializing central powers. The next
argument deals with those possible cancellations explicitly.

\begin{lemma}[Residue coordinates]\label{lem:coords}
For $i=5a+r$ and $j=5b+s$, with $a,b\in\ZZ$ and $0\leq r,s<5$,
\begin{equation}\label{eq:residue}
U^iV^j=x^ay^bU^rV^s.
\end{equation}
For the numerator \eqref{eq:Nexpanded}, the coordinates on $UV^3$ and
$U^2V^3$ in the basis \eqref{eq:basis} are respectively
\begin{align}
C_{13}&=\frac{q(x^2y+xy^2-2)}{x^2y^2},\label{eq:C13}\\
C_{23}&=-\frac{q^3(2x^2y-xy^2-1)}{xy^2}.\label{eq:C23}
\end{align}
\end{lemma}
\begin{proof}
The equality $U^{5a+r}=x^aU^r$ is valid for positive and negative $a$,
and similarly for $V$. The extracted powers of $x,y$ are central,
so moving them to the left introduces no phase. This proves
\eqref{eq:residue}.

Exactly three terms of Appendix~\ref{app:N} have residues $(1,3)$:
\[
(-9,-7;-2q),\qquad(-4,3;q),\qquad(1,-2;q).
\]
Their combined coefficient is
$-2qx^{-2}y^{-2}+qx^{-1}+qy^{-1}$, which becomes
\eqref{eq:C13} after a common denominator is taken. Exactly three terms
have residues $(2,3)$:
\[
(-3,-7;q^3),\qquad(2,3;q^3),\qquad(7,-2;-2q^3).
\]
Their coefficient is $q^3x^{-1}y^{-2}+q^3-2q^3xy^{-1}$, giving
\eqref{eq:C23}. The remaining table entries belong to other residue
classes and cannot contribute to these coordinates.
\end{proof}

\begin{proposition}[Nonscalarity in every division specialization]\label{prop:nonscalar}
One has $N\notin K$. In particular $N\ne0$ and, as an element of $D$,
$N\notin F$.
\end{proposition}
\begin{proof}
The basis vector indexed by $(0,0)$ is $1$. Therefore an element of $K$
has zero coordinates on all other basis vectors. If $N\in K$, both
\eqref{eq:C13} and \eqref{eq:C23} vanish. Their denominators and $q$ are
nonzero, so, on putting $A_0=x^2y$ and $B_0=xy^2$, we obtain
\[
A_0+B_0=2,\qquad 2A_0-B_0=1.
\]
Adding the two equations gives $3A_0=3$. Characteristic zero implies
$A_0=1$, and the first equation then gives $B_0=1$. Hence
$0=A_0-B_0=xy(x-y)$. Since $xy\ne0$, this forces $x=y$, contradicting
Lemma~\ref{lem:distinct}. Thus $N\notin K$. Since $0\in K$ and
$F\subseteq K$, the final assertions follow.
\end{proof}

The proof requires an actual basis, not merely a spanning list of $25$
monomials. It also explains why a proof using a single generic coefficient
would be insufficient for the stated family. The two-coordinate
contradiction is insensitive to algebraic relations among $x,y$, provided
the given algebra is a division algebra with the required basis.

\section{The exact fifth power certificate}\label{sec:certificate}

\subsection{An integral twisted Laurent algebra}

Let
\[
R=\ZZ[Q]/(Q^4+Q^3+Q^2+Q+1),
\]
and write $Q$ for the residue class. It satisfies $Q^5=1$ and is a unit.
Every element of $R$ is represented uniquely by an integer quadruple
$(a_0,a_1,a_2,a_3)$, since division by the monic polynomial gives a unique
remainder of degree at most three. Introduce a free $R$-module with basis
$e_{ij}$ indexed by $(i,j)\in\ZZ^2$ and finite support. Define
\begin{equation}\label{eq:twisted}
e_{ij}e_{kl}=Q^{-jk}e_{i+k,j+l}.
\end{equation}
This multiplication is associative. On multiplying three basis vectors,
the exponents of $Q$ from the two parenthesizations are
\[
-jk-(j+l)m\quad\hbox{and}\quad-lm-j(k+m),
\]
which are equal. Bilinearity extends associativity to finite sums.
The unit is $e_{00}$.

Equation~\eqref{eq:cyclotomic} follows from $q^5=1$ and $q\ne1$:
$(q-1)(1+q+q^2+q^3+q^4)=q^5-1=0$, so the second factor vanishes in
$K$. Thus $Q\mapsto q$ defines a ring map $R\to K$. By
Lemma~\ref{lem:monomial}, the map
\begin{equation}\label{eq:evaluation}
\ev(e_{ij})=U^iV^j
\end{equation}
extends to a ring homomorphism from the twisted Laurent algebra to $A$.
It need not be injective: relations among $x,y$ may identify central
Laurent expressions. Injectivity is unnecessary for transporting a
universal identity. The actual symbol basis is used only when extracting
coordinates for nonscalarity and dimension.

Write $\mathcal N=\sum n_{ij}(Q)e_{ij}$ for the formal table in
Appendix~\ref{app:N}. Define $H(X,Y)$ by Appendix~\ref{app:H}:
\begin{equation}\label{eq:H}
H(X,Y)=\sum_{(a,b)}h_{ab}X^aY^b\in\ZZ[X^{\pm1},Y^{\pm1}].
\end{equation}
There are $75$ nonzero entries in that table, all ordinary integers.

\begin{proposition}[The exact fifth-power identity]\label{prop:fifth}
In the twisted Laurent algebra one has
\begin{equation}\label{eq:universal-fifth}
\mathcal N^5=\sum_{(a,b)}h_{ab}e_{5a,5b}.
\end{equation}
Thus every coefficient at an exponent outside $5\ZZ^2$ is zero, and the
right-hand side is exactly the table in Appendix~\ref{app:H}.
\end{proposition}
\begin{proof}[Exact finite certificate]
Apply the coefficient reduction and convolution specified below to the
$69$-term table. The complete products have respectively $237$, $552$,
$969$, and $75$ nonzero terms. Every multiplication equality is established
by finite exact arithmetic; the last output is coefficientwise equal to
the $75$-term table, with exponents multiplied by five. These equalities,
not merely the support counts, prove \eqref{eq:universal-fifth} using
\eqref{eq:twisted}.

The checked proof package supplies all multiplication certificates in the
full fifth-power chain, a soundness theorem transporting their operations
to evaluated products, and a kernel-checked equality identifying the final
literal output with the standalone $75$-term candidate list. The checked
support theorem and integer-coefficient theorem apply to that same list.
The complete evaluated theorem is
\code{PrimeCyclicity.N5Modern.numerator_fifth}; final alignment is
\code{PrimeFiveFamily.fifthOutput_eq_candidate}. The latter has no axioms.
The exact recurrence below and the complete input and output tables allow
independent arithmetic replay. No assertion of equality is inferred solely
from a central-looking output list or a numerical specialization.
\end{proof}

\paragraph{Status of the certificate.}
The first full Lean check is complete, including semantic evaluation and
final-family integration. Recorded external exact-arithmetic computations
and the standalone script supplied with this paper also reproduce the
identity. The independent N4 and structural rebuilds, fresh N5 replay,
and downstream final-assembly checks have passed. Original earlier
base/N2/N3 objects are reused, so this is not a fresh rebuild of the entire
fifth-power dependency closure. The Lean export is the universal
evaluated identity in its specified algebraic interface. The twisted
Laurent algebra here is the mathematical presentation of the same finite
arithmetic, not a claim that Lean separately implements this particular
quotient-and-free-module definition.

\subsection{The finite computation in reproducible form}

For clarity, the full arithmetic algorithm is specified independently of
a particular programming language. Represent a Laurent expression as a
finite map $C\colon\ZZ^2\to\ZZ^4$, with omitted entries interpreted as zero.
For coefficient quadruples $a,b$ and an integer $s$, let
\[
\operatorname{cmul}(a,b;s)
 =\red\bigl(Q^{\bar s}(a_0+a_1Q+a_2Q^2+a_3Q^3)
                     (b_0+b_1Q+b_2Q^2+b_3Q^3)\bigr),
\]
where $\bar s\in\{0,1,2,3,4\}$ is congruent to $s$ modulo five and
$\red$ is the monic remainder modulo $1+Q+Q^2+Q^3+Q^4$.
To compute the latter, process the highest nonzero degree $d\geq4$,
remove its coefficient $tQ^d$, and subtract $t$ from the coefficients
at degrees $d-4,d-3,d-2,d-1$. This replaces
$Q^d$ by $-Q^{d-4}-Q^{d-3}-Q^{d-2}-Q^{d-1}$ and strictly lowers the
largest degree still requiring reduction. It terminates and gives the
unique remainder.

The product of maps $C,E$ is
\begin{equation}\label{eq:convolution}
(C\ast E)_{r,s}
 =\sum_{\substack{i+k=r\\j+l=s}}
     \operatorname{cmul}(C_{i,j},E_{k,l};-jk).
\end{equation}
Only pairs in the two finite supports enter the sum. Add quadruples
componentwise and delete entries that become the zero quadruple. Starting
with $C^{(1)}=(n_{ij})$, form
$C^{(m+1)}=C^{(m)}\ast C^{(1)}$ for $m=1,2,3,4$.
An exhaustive check compares $C^{(5)}$ with the right-hand map in
\eqref{eq:universal-fifth}. The comparison must be equality of finite maps,
not agreement on a selection of positions: otherwise noncentral residual
terms could be missed.

The recorded nonzero support sizes are
\[
69,\quad237,\quad552,\quad969,\quad75
\]
for powers one through five. The first four maps have no support in
$5\ZZ^2$; all support of the fifth map is in $5\ZZ^2$. Those counts are
useful diagnostics but do not prove equality. Many different maps have
the same support size. Only coefficientwise equality, with all omitted
coefficients checked to be zero, supplies the identity.

\subsection{Why the arithmetic gives centrality after specialization}

\begin{lemma}[Semantic consequence of the certificate]\label{lem:fifth-semantic}
One has
\[
N^5=H(x,y)\in K.
\]
Moreover $H(x,y)\in F^\times$.
\end{lemma}
\begin{proof}
Apply the ring homomorphism \eqref{eq:evaluation} to
\eqref{eq:universal-fifth}. The left side becomes $N^5$, not merely the
evaluation of an unrelated output list. For each right-hand term,
$U^{5a}V^{5b}=x^ay^b$ by \eqref{eq:residue}; the result is $H(x,y)$.

Proposition~\ref{prop:bridge} gives $\sigma(N)=N$. Hence
\[
\tau(H(x,y))=\sigma(H(x,y))=\sigma(N^5)=\sigma(N)^5=N^5=H(x,y).
\]
The equality is interpreted in $K$ using its injective embedding in $A$,
so $H(x,y)\in F$. By Proposition~\ref{prop:nonscalar}, $N\ne0$.
A nonzero element of a division ring has nonzero fifth power. Therefore
$H(x,y)\ne0$.
\end{proof}

This last nonvanishing conclusion does not require a separate analysis of
the zeros of the $75$-term polynomial. It is a consequence of the exact
identity and uniform nonscalarity in the given division algebra.

\subsection{Arithmetic certificates versus a trusted computation}

A generator may compute candidate lists and place them in Lean source,
but that generator is not a logical oracle. The checking layer proves
finite multiplication equalities by kernel reduction of exact integer
and list operations. Splitting a product into rows or smaller pieces
changes the proof's size, not the required equality. Each row has to be
joined to the exact partial sum from the preceding rows, and the exported
theorem has to cover the whole input list.

A second layer proves that the finite operations agree with evaluation in
an arbitrary division algebra satisfying the scalar and commutation
hypotheses. A closed finite-list certificate alone does not prove that its
output evaluates to $N^5$. Conversely, a generic evaluation-multiplication
lemma alone does not establish that the specific output has central support.
The final result needs both layers and the link to the frozen numerator.

For the completed cube stage, the recorded audit reports no axioms for the
closed arithmetic certificates and only the standard logical principles
\code{propext}, \code{Classical.choice}, and \code{Quot.sound} for the
semantic theorem. These reports concern their named declarations and
imports. The separate full fifth-power and final-assembly audits now supply the
corresponding checked coverage for those completed stages. Hashes identify a snapshot; they do not independently
prove a mathematical statement.

\section{The cyclic maximal subfield endpoint}\label{sec:endpoint}

We give the field argument explicitly. In this section, ``cyclic division
algebra'' means a central division algebra containing a maximal subfield
that is a cyclic Galois extension of its center. This is the usual
maximal-subfield characterization; no presentation by a second cyclic
algebra generator is needed for the conclusion stated here.

\begin{lemma}[A fifth-central element in central dimension twenty five]
\label{lem:endpoint}
Let $B$ be a division algebra with center $E$, $\dim_E B=25$, and
$\operatorname{char}E=0$. Suppose $q\in E$ has order five and $r\in B$
satisfies $r\notin E$ and $r^5=h\in E$. Then $E[r]=E(r)$ is a cyclic
Galois extension of $E$ of degree five and is a maximal commutative
subfield of $B$.
\end{lemma}
\begin{proof}
Because $E$ is central, the subring $E[r]$ is commutative. The equation
$r^5=h$ shows that it is spanned over $E$ by $1,r,r^2,r^3,r^4$.
It is a domain, being a subring of a division ring. Multiplication by any
nonzero element $s\in E[r]$ is an injective linear endomorphism of this
finite-dimensional $E$-space and hence is surjective. In particular some
$t\in E[r]$ satisfies $st=1$. Thus $L=E[r]$ is a field, and equals $E(r)$.

Write $d=[L:E]$. The preceding spanning argument gives $d\leq5$.
Regard $B$ as a left vector space over $L$. The scalar action need not be
central for this left vector-space structure: multiplication in $B$ is
associative. A finite $L$-basis exists because an $E$-spanning set is also
an $L$-spanning set. If $\ell_1,\ldots,\ell_d$ is an $E$-basis of $L$
and $b_1,\ldots,b_m$ is a left $L$-basis of $B$, the products
$\ell_i b_j$ form an $E$-basis of $B$. Spanning follows by expanding each
left coefficient, and independence follows first from the $L$-independence
of the $b_j$ and then from the $E$-independence of the $\ell_i$.
Consequently $25=dm$, so $d$ divides $25$. Since $r\notin E$, $d>1$;
together with $d\leq5$ this forces $d=5$.

The monic minimal polynomial of $r$ over $E$ divides $X^5-h$ and has
its full degree five. Hence it equals $X^5-h$. Also $r\ne0$, since
$0\in E$, so $h\ne0$. The five elements $q^i r$ belong to $L$, are
pairwise distinct, and each has fifth power $h$. Thus $X^5-h$ splits
into distinct linear factors in $L[X]$. Its irreducibility and the fact
that $L$ is its splitting field show that $L/E$ is Galois. More explicitly,
substitution $r\mapsto qr$ induces an $E$-automorphism of
$E[X]/(X^5-h)\simeq L$: it respects the relation and has inverse
$r\mapsto q^{-1}r$. Its successive powers send $r$ to $q^i r$ and have
order exactly five. Since a degree-five extension has at most five
$E$-automorphisms, this automorphism generates $\Gal(L/E)$.

Finally, let $M$ be any commutative $E$-subalgebra of $B$ containing $L$.
It is finite dimensional and a domain, so the same multiplication argument
makes it a field. The vector-space argument just given shows $[M:E]$
divides $25$, and the field tower through $L$ shows it is a multiple of
five. Therefore $[M:E]$ is $5$ or $25$. In the former case $M=L$.
In the latter case the $E$-subspace $M\subseteq B$ has the full dimension
of $B$, hence $M=B$; this would make $B$ commutative, contradicting
$\Cent(B)=E$ and $\dim_E B=25>1$. Thus $M=L$, proving maximality.
\end{proof}

\begin{theorem}[Explicit cyclicity for the specified family]
\label{thm:main}
Let the data satisfy Definition~\ref{def:data}, and let $N$ be the exact
numerator of Appendix~\ref{app:N}, equivalently \eqref{eq:Ncompact}.
Then $D/F$ is
a central division algebra of degree five, and
\[
L=F(N)\subseteq D
\]
is a cyclic maximal subfield of degree five. Its minimal polynomial is
$X^5-H(x,y)$ and a generator of $\Gal(L/F)$ sends $N$ to $qN$.
\end{theorem}
\begin{proof}
Theorem~\ref{thm:descent} gives the central dimension and center of $D$.
Propositions~\ref{prop:bridge} and \ref{prop:nonscalar} put $N$ in $D$
but outside $F$. Proposition~\ref{prop:fifth} and Lemma~\ref{lem:fifth-semantic} give
$N^5=H(x,y)\in F^\times$. Lemma~\ref{lem:cycle} places a primitive fifth
root $q$ in $F$. Apply Lemma~\ref{lem:endpoint} with
$(B,E,r,h)=(D,F,N,H(x,y))$.
\end{proof}

\paragraph{Scope of the checked final export.}
The final export states degree five, Galoisness, cyclicity of the Galois
group, and maximality of the generated commutative subalgebra, with no
additional scalar fifth-power hypothesis. It obtains the required identity
from the completed power certificate and applies the earlier conditional
endpoint internally. The identification of the minimal polynomial as
$X^5-H(x,y)$ and the designation of the generator $N\mapsto qN$ are proved
in the paper argument above. They are not additional conclusions directly
named in that final export. The relative-family structural assumptions
remain essential; removing the intermediate power hypothesis does not
remove any of those assumptions.

\section{Relationship with the literature and limits of the claim}
\label{sec:literature}

The motivation is an existing fixed-ring construction. Fryer studies the
order-three monomial action on a quantum division ring and constructs
fixed quantum generators under an auxiliary primitive cube-root
hypothesis \cite[Proposition~3.3 and Theorem~3.4]{Fryer}. At a fifth root
$q$, roots of $q$ needed there can be chosen among powers of $q$, but a
primitive cube root of unity need not be in the ground field. The present
construction addresses a cyclicity conclusion for the restricted data in
Definition~\ref{def:data}, including division-algebra specializations;
it does not reproduce a full fixed-ring isomorphism theorem. Fryer's
Questions~4.12 do not explicitly formulate removal of the cube-root
hypothesis as the problem treated here. No such open-problem attribution
is made.

Reichstein and Vonessen discuss the related monomial action within an
action of $(C_p)^2\rtimes C_3$ on a generic symbol algebra
\cite[Section~15, especially Proposition~15.3]{ReichsteinVonessen}.
Their obstruction concerns a maximal subfield invariant under the entire
larger group. The present construction seeks a field in the $C_3$-fixed
algebra, with no requirement of invariance under the additional $(C_5)^2$.
These are different conditions. On exponent pairs the $C_3$ action has
matrix
\[
\begin{pmatrix}0&-1\\1&-1\end{pmatrix},
\]
whose characteristic polynomial $T^2+T+1$ has no root in $\FF_5$.
This simple calculation explains why an invariant monomial line is not the
natural source of the nonmonomial candidate $N$.

Matzri's theorem proves cyclicity for degree-five division algebras with
a specified dihedral splitting field
\cite[Theorem~3.2]{Matzri}. The order-three semilinear action in our data
does not supply that hypothesis. In particular, an argument that cyclicity
holds after adjoining a cube root of unity cannot be replaced by an
unjustified assertion that cyclicity descends through every quadratic
extension. The distinction between cyclicity questions and particular
splitting-field structures is also emphasized in the survey
\cite[Section~2]{ABGV}. That survey supplies historical context, not a
claim about the present status of every unrestricted cyclicity problem.

A bounded targeted review of primary literature did not identify the exact
root-free specialization-uniform statement or the exact numerator presented
here. This negative search result does not establish originality. The
family theorem, the explicit invariant, and its machine-checkable
certificate would need separate comparison with the literature before any
claim of novelty or priority. In particular, this paper is not described
as a solution to an explicitly published cube-root-removal question.

\section{Formal coverage and verification limits}\label{sec:formal}

The mathematical data correspond to
\code{PrimeFiveFamily.SymbolFamily}. Its fields specify a characteristic-zero
scalar field, a division algebra, a primitive fifth root, actual units,
the quantum relation, scalar fifth powers, the actual monomial basis,
the center characterization, and the two compatible automorphisms.
The structure does not contain the fixed dimension, the fixed center,
cyclicity, or any assertion about the numerator as input fields.

\subsection{Structural and numerator interfaces}

The structural family manifest records $26$ audited declarations. An
independent structural rebuild checked all $34$ modules in the structural
dependency chain from fresh project objects, with zero reused project
objects. Official upstream Mathlib and dependency objects were reused;
the receipt does not claim those libraries were rebuilt. This earlier
rebuild ended at the explicitly conditional supporting endpoint. It is
separate from the subsequent complete power and final-assembly checks.

The named interfaces make the connection to the mathematical argument
precise.
\begin{itemize}[leftmargin=1.6em]
\item \code{frozen_eq_compact} identifies the evaluated frozen $69$-term
list with the compact expression, accounting for the actual $U,V$ and
commutation phase.
\item \code{frozenNumerator_fixed} and
\code{frozenNumerator_not_scalar} establish invariance and nonscalarity
of that exact list. The definition \code{fixedNumerator} places it in the
actual fixed algebra.
\item \code{tau_x}, \code{tau_y}, \code{tau_z}, and \code{x_ne_y}
derive the scalar cycle and inequality. The names
\code{fixed_finrank}, \code{fixedDivisionRing}, and \code{fixed_center}
record the elementary descent consequences.
\item The earlier \code{supporting_cyclicity_of_fifthPower_identity}
retains an explicit identity hypothesis. It remains a useful conditional
lemma, but the final assembly now supplies that identity as a theorem.
\end{itemize}

\subsection{Completed fifth power and final assembly}

The complete evaluated fifth-power export is
\code{PrimeCyclicity.N5Modern.numerator_fifth}, in
\code{modern.N5Fifth}. Its literal output is
\code{PrimeCyclicity.N5Modern.acc69}, with $75$ terms. The recorded
fifth-power audit covers $416$ chain modules and reports matching source
and main-object hashes. Closed arithmetic certificates have no axioms;
the semantic theorem reports only \code{propext},
\code{Classical.choice}, and \code{Quot.sound}.

The separate checked helper
\code{FrozenCentralOutput75.candidateOutput75_support} proves that every
exponent in the expected output is divisible by five; its audit reports
only \code{propext}. Its length and integer-coefficient assertions have
no axioms. Crucially, final assembly proves
\code{fifthOutput_eq_candidate}, identifying the completed literal output
with that exact helper list by kernel computation and without axioms.
The support certificate therefore applies to the actual power output,
not merely a proposed table.

Next, \code{frozenNumerator_fifth} instantiates the universal evaluated
identity with the actual scalar unit $q$ and actual units $U,V$, then
applies the central-support evaluator using $U^5=x$ and $V^5=y$.
The scalar is defined by
\code{fifthScalar}, the central coefficient evaluation of the matched
$75$-term list. Its mathematical value is $H(x,y)$ by the exponent
normalization in Appendix~\ref{app:H}.

The final theorem is
\begin{center}
\code{PrimeFiveFamily.symbolFamily_cyclicity}.
\end{center}
It proves degree five, Galoisness, cyclicity of the Galois group, and
maximality among commutative subalgebras for the field generated by the
actual fixed numerator. Its only data are the specified
\code{SymbolFamily} and ambient field, characteristic-zero, division-ring,
and algebra instances. There is no extra power, fixedness, nonscalarity,
dimension, or cyclicity hypothesis. The final assembly and audit passed
Lean 4.32.2 against the pinned Mathlib 4.32.2 environment using existing
dependencies. Their final axiom report is exactly
\code{propext}, \code{Classical.choice}, and \code{Quot.sound}.

The initial check and a staged replay of the relative-family theorem
have passed. The final replay freshly compiled $418$ N5 modules, a
separate N5 audit, and \code{FrozenCentralOutput75},
\code{FinalAssembly}, and \code{FinalAudit}: $422$ successful
invocations, with no failures or retries. It reused $212$ independently
rebuilt N4 modules, $31$ independently rebuilt structural modules, and
$99$ original earlier base/N2/N3 module objects. Of those $99$,
\code{modern.Algebra} and \code{modern.Numerator} also had separate
source rebuilds; the remaining $97$ retain original checked-object
provenance. Source and object hashes were verified. This is not an
independent fresh rebuild of all $763$ custom sources. The included
public verifier has not completed an end-to-end replay. Source files,
receipts, statements, and hashes identify the checked artifact; a hash
itself is not a proof.

\subsection{Replay evidence and review limits}

The terminal replay used $3348.363222$ child CPU seconds. Its exact
module inventory, per-source hashes, compiler receipts, final theorem
statement, and axiom report accompany the source package. The replay
used pinned upstream compiled objects; no whole-Mathlib source rebuild,
separate-kernel implementation, fresh-network installation test, or
external human peer review is claimed.

A bounded literature search does not establish originality. A specialist
assessment would be needed before making a novelty or priority claim.
The minimal-polynomial and designated-generator claims are mathematical
arguments in this paper, distinct from the final Lean export. No author
or runtime model identity is asserted by the paper.

The scope remains Definition~\ref{def:data}. No generic Ore division ring
has been constructed by the formal theorem, no positive-characteristic
extension is asserted, and no unrestricted degree-five or prime-degree
cyclicity problem is solved.

\appendix
\section{The complete numerator table}\label{app:N}

The following table specifies $N$ without reliance on a hidden coefficient
file. Each row $(i,j;a_0,a_1,a_2,a_3)$ contributes
$(a_0+a_1q+a_2q^2+a_3q^3)U^iV^j$ to \eqref{eq:Nexpanded}. All omitted
coefficients are zero. The order of $U$ and $V$ is part of the definition.
\begin{longtable}{rrrrrr}
\toprule
$i$ & $j$ & $a_0$ & $a_1$ & $a_2$ & $a_3$ \\
\midrule
\endfirsthead
\toprule
$i$ & $j$ & $a_0$ & $a_1$ & $a_2$ & $a_3$ \\
\midrule
\endhead
\midrule
\multicolumn{6}{r}{Continued on the next page} \\
\endfoot
\bottomrule
\endlastfoot
$-10$ & $-9$ & $0$ & $-1$ & $0$ & $0$ \\
$-9$ & $-10$ & $-1$ & $0$ & $0$ & $0$ \\
$-9$ & $-7$ & $0$ & $-2$ & $0$ & $0$ \\
$-8$ & $-8$ & $-1$ & $-2$ & $-2$ & $-1$ \\
$-8$ & $-5$ & $1$ & $1$ & $1$ & $1$ \\
$-7$ & $-9$ & $0$ & $0$ & $0$ & $-2$ \\
$-7$ & $-6$ & $0$ & $-2$ & $-1$ & $-1$ \\
$-7$ & $-3$ & $1$ & $0$ & $0$ & $0$ \\
$-6$ & $-7$ & $-1$ & $1$ & $1$ & $0$ \\
$-6$ & $-4$ & $-1$ & $-1$ & $1$ & $0$ \\
$-6$ & $-1$ & $-2$ & $-2$ & $-2$ & $-2$ \\
$-5$ & $-8$ & $1$ & $1$ & $1$ & $1$ \\
$-5$ & $-2$ & $-2$ & $-3$ & $-1$ & $-1$ \\
$-5$ & $1$ & $0$ & $2$ & $0$ & $0$ \\
$-4$ & $-6$ & $0$ & $1$ & $1$ & $2$ \\
$-4$ & $-3$ & $-2$ & $-2$ & $-1$ & $0$ \\
$-4$ & $0$ & $-1$ & $-1$ & $0$ & $0$ \\
$-4$ & $3$ & $0$ & $1$ & $0$ & $0$ \\
$-3$ & $-7$ & $0$ & $0$ & $0$ & $1$ \\
$-3$ & $-4$ & $0$ & $-1$ & $-2$ & $-2$ \\
$-3$ & $-1$ & $0$ & $0$ & $0$ & $1$ \\
$-3$ & $2$ & $2$ & $1$ & $1$ & $-1$ \\
$-3$ & $5$ & $1$ & $1$ & $1$ & $1$ \\
$-2$ & $-5$ & $1$ & $3$ & $2$ & $2$ \\
$-2$ & $-2$ & $0$ & $2$ & $0$ & $0$ \\
$-2$ & $1$ & $0$ & $0$ & $0$ & $1$ \\
$-2$ & $4$ & $0$ & $1$ & $-1$ & $-1$ \\
$-2$ & $7$ & $-2$ & $0$ & $0$ & $0$ \\
$-1$ & $-6$ & $2$ & $0$ & $0$ & $0$ \\
$-1$ & $-3$ & $1$ & $0$ & $0$ & $0$ \\
$-1$ & $3$ & $-1$ & $-2$ & $-2$ & $0$ \\
$-1$ & $6$ & $-1$ & $-1$ & $-2$ & $0$ \\
$-1$ & $9$ & $1$ & $1$ & $1$ & $1$ \\
$0$ & $-4$ & $-1$ & $-1$ & $0$ & $0$ \\
$0$ & $2$ & $-2$ & $-2$ & $-2$ & $-2$ \\
$0$ & $8$ & $1$ & $0$ & $-1$ & $-1$ \\
$1$ & $-5$ & $2$ & $0$ & $0$ & $0$ \\
$1$ & $-2$ & $0$ & $1$ & $0$ & $0$ \\
$1$ & $4$ & $0$ & $-2$ & $-2$ & $-1$ \\
$1$ & $7$ & $1$ & $1$ & $-1$ & $0$ \\
$1$ & $10$ & $0$ & $-1$ & $0$ & $0$ \\
$2$ & $-3$ & $-1$ & $1$ & $1$ & $2$ \\
$2$ & $0$ & $-2$ & $-2$ & $-2$ & $-2$ \\
$2$ & $3$ & $0$ & $0$ & $0$ & $1$ \\
$2$ & $6$ & $0$ & $2$ & $1$ & $1$ \\
$2$ & $9$ & $0$ & $0$ & $0$ & $-2$ \\
$3$ & $-4$ & $0$ & $0$ & $0$ & $1$ \\
$3$ & $-1$ & $0$ & $1$ & $2$ & $2$ \\
$3$ & $2$ & $1$ & $0$ & $0$ & $0$ \\
$3$ & $5$ & $-2$ & $-3$ & $-1$ & $-1$ \\
$3$ & $8$ & $0$ & $-1$ & $0$ & $0$ \\
$4$ & $-2$ & $2$ & $1$ & $1$ & $0$ \\
$4$ & $1$ & $2$ & $2$ & $1$ & $0$ \\
$4$ & $4$ & $0$ & $1$ & $1$ & $1$ \\
$4$ & $7$ & $1$ & $0$ & $0$ & $0$ \\
$5$ & $-3$ & $1$ & $1$ & $1$ & $1$ \\
$5$ & $3$ & $1$ & $3$ & $2$ & $2$ \\
$5$ & $6$ & $2$ & $0$ & $0$ & $0$ \\
$6$ & $-1$ & $0$ & $1$ & $1$ & $-1$ \\
$6$ & $2$ & $1$ & $1$ & $2$ & $0$ \\
$6$ & $5$ & $0$ & $2$ & $0$ & $0$ \\
$7$ & $-2$ & $0$ & $0$ & $0$ & $-2$ \\
$7$ & $1$ & $0$ & $-1$ & $1$ & $1$ \\
$7$ & $4$ & $0$ & $0$ & $0$ & $1$ \\
$8$ & $0$ & $1$ & $0$ & $-1$ & $-1$ \\
$8$ & $3$ & $0$ & $-1$ & $0$ & $0$ \\
$9$ & $-1$ & $-1$ & $0$ & $0$ & $0$ \\
$9$ & $2$ & $-2$ & $0$ & $0$ & $0$ \\
$10$ & $1$ & $-1$ & $0$ & $0$ & $0$ \\
\end{longtable}


\section{The central fifth power}\label{app:H}

Each row $(a,b;h)$ contributes $hX^aY^b$ to $H(X,Y)$. The exponent pairs
here are divided by five relative to the output Laurent list in
$U,V$. Thus this table has $75$ terms in $X,Y$ and
\eqref{eq:universal-fifth} uses $(5a,5b)$ in the twisted Laurent algebra.
All omitted coefficients are zero. Proposition~\ref{prop:fifth} and the
checked literal-output alignment establish that this is the actual fifth
power, not merely a central-looking candidate table.
\begin{longtable}{rrr}
\toprule
$a$ & $b$ & $h_{ab}$ \\
\midrule
\endfirsthead
\toprule
$a$ & $b$ & $h_{ab}$ \\
\midrule
\endhead
\midrule
\multicolumn{3}{r}{Continued on the next page} \\
\endfoot
\bottomrule
\endlastfoot
$-10$ & $-9$ & $-1$ \\
$-9$ & $-10$ & $-1$ \\
$-9$ & $-7$ & $3$ \\
$-8$ & $-8$ & $14$ \\
$-8$ & $-5$ & $-6$ \\
$-7$ & $-9$ & $3$ \\
$-7$ & $-6$ & $-39$ \\
$-7$ & $-3$ & $6$ \\
$-6$ & $-7$ & $-39$ \\
$-6$ & $-4$ & $84$ \\
$-6$ & $-1$ & $-3$ \\
$-5$ & $-8$ & $-6$ \\
$-5$ & $-5$ & $60$ \\
$-5$ & $-2$ & $-117$ \\
$-5$ & $1$ & $-3$ \\
$-4$ & $-6$ & $84$ \\
$-4$ & $-3$ & $-90$ \\
$-4$ & $0$ & $138$ \\
$-4$ & $3$ & $6$ \\
$-3$ & $-7$ & $6$ \\
$-3$ & $-4$ & $-90$ \\
$-3$ & $-1$ & $21$ \\
$-3$ & $2$ & $-117$ \\
$-3$ & $5$ & $-6$ \\
$-2$ & $-5$ & $-117$ \\
$-2$ & $-2$ & $192$ \\
$-2$ & $1$ & $21$ \\
$-2$ & $4$ & $84$ \\
$-2$ & $7$ & $3$ \\
$-1$ & $-6$ & $-3$ \\
$-1$ & $-3$ & $21$ \\
$-1$ & $0$ & $-120$ \\
$-1$ & $3$ & $-90$ \\
$-1$ & $6$ & $-39$ \\
$-1$ & $9$ & $-1$ \\
$0$ & $-4$ & $138$ \\
$0$ & $-1$ & $-120$ \\
$0$ & $2$ & $192$ \\
$0$ & $5$ & $60$ \\
$0$ & $8$ & $14$ \\
$1$ & $-5$ & $-3$ \\
$1$ & $-2$ & $21$ \\
$1$ & $1$ & $-120$ \\
$1$ & $4$ & $-90$ \\
$1$ & $7$ & $-39$ \\
$1$ & $10$ & $-1$ \\
$2$ & $-3$ & $-117$ \\
$2$ & $0$ & $192$ \\
$2$ & $3$ & $21$ \\
$2$ & $6$ & $84$ \\
$2$ & $9$ & $3$ \\
$3$ & $-4$ & $6$ \\
$3$ & $-1$ & $-90$ \\
$3$ & $2$ & $21$ \\
$3$ & $5$ & $-117$ \\
$3$ & $8$ & $-6$ \\
$4$ & $-2$ & $84$ \\
$4$ & $1$ & $-90$ \\
$4$ & $4$ & $138$ \\
$4$ & $7$ & $6$ \\
$5$ & $-3$ & $-6$ \\
$5$ & $0$ & $60$ \\
$5$ & $3$ & $-117$ \\
$5$ & $6$ & $-3$ \\
$6$ & $-1$ & $-39$ \\
$6$ & $2$ & $84$ \\
$6$ & $5$ & $-3$ \\
$7$ & $-2$ & $3$ \\
$7$ & $1$ & $-39$ \\
$7$ & $4$ & $6$ \\
$8$ & $0$ & $14$ \\
$8$ & $3$ & $-6$ \\
$9$ & $-1$ & $-1$ \\
$9$ & $2$ & $3$ \\
$10$ & $1$ & $-1$ \\
\end{longtable}


\section{Reproduction and provenance}\label{app:reproduce}

The companion file \code{certificate-data.json} contains exactly the two
tables, with integer coefficient quadruples for the numerator and integer
coefficients for $H$. The file \code{replay-certificate.py} uses only the
Python standard library. It implements monic coefficient reduction,
twisted Laurent multiplication, the pair construction, exhaustive
fifth-power comparison, the two residue-coordinate checks, and the
order-three action on Laurent monomials. It is a review aid and external
exact-arithmetic check, not a substitute for the Lean kernel.

The companion \code{source-snapshot.json} records the aggregate source
archive hash, source count, final theorem name, and snapshot verification
status. It points to the per-source checksum manifest; it does not contain
stage-by-stage build receipts. The package verification record and
Section~\ref{sec:formal} distinguish the initial checks, the staged replay
status, reused objects, and review caveats.

\bibliographystyle{alpha}
\bibliography{references}
\end{document}
